% !TeX program = pdflatex
% Standalone IEEE Transactions experiment section for IRPP-TD.
% RQ1--RQ3 are backfilled with verified rerun results; RQ4--RQ5 remain placeholders.
\documentclass[journal]{IEEEtran}

\usepackage[T1]{fontenc}
\usepackage[utf8]{inputenc}
\usepackage{textcomp}
\usepackage{amsmath,amsfonts,amssymb}
\usepackage{array}
\usepackage{booktabs}
\usepackage{graphicx}
\usepackage{multirow}
\usepackage{tabularx}
\usepackage{makecell}
\usepackage{xcolor}
\usepackage{stfloats}
\usepackage{placeins}

\definecolor{tbdblue}{RGB}{23,74,135}
\definecolor{panelgray}{RGB}{246,247,249}
\newcolumntype{Y}{>{\centering\arraybackslash}X}
\newcommand{\tbd}{\textcolor{tbdblue}{\textsf{TBD}}}
\newcommand{\fillnote}[1]{\textcolor{tbdblue}{\textit{[#1]}}}
\newcommand{\panelbox}[2]{%
  \fcolorbox{black!45}{panelgray}{%
    \parbox[c][0.66in][c]{0.222\textwidth}{%
      \centering\scriptsize\textbf{#1}\par\vspace{0.35ex}%
      \footnotesize #2}}}

\hyphenation{op-tical net-works semi-conduc-tor block-chain crowd-sens-ing}

\begin{document}
\markboth{IEEE Transactions on Mobile Computing, RQ1--RQ3 Experiment Revision}%
{IRPP-TD: Performance Evaluation}
\setcounter{section}{6}

\section{Performance Evaluation}
\label{sec:performance-evaluation}

\subsection{Research Questions and Methodology}
\label{sec:experimental-methodology}

We evaluate IRPP-TD within the threat-model boundary. The intended data
requester (DR) decrypts reports and observes task participation and reputation;
accordingly, the experiments assess TD quality, worker-side robustness,
efficiency, and conditional accountability, rather than confidentiality
against the DR. We organize the evaluation around five questions.

\begin{itemize}
  \setlength{\itemsep}{0.15ex}
  \setlength{\topsep}{0.3ex}
  \setlength{\parsep}{0pt}
  \item \emph{RQ1---Effectiveness and applicability:} How accurately and
  consistently does IRPP-TD recover truth across heterogeneous task-aligned
  sensing reports, workloads, task phases, and participation levels?
  \item \emph{RQ2---Cold-start efficiency and recovery:} Can IRPP-TD resolve
  cold start with few trusted reports and match exact-ABOD accuracy at lower
  angular cost? Which components enable recovery from early errors?
  \item \emph{RQ3---Worker-side poisoning:} How accurately and reliably does
  IRPP-TD operate as malicious prevalence and normalized attack strength grow
  under independent, compact-collusive, on--off, and mature-anchor poisoning,
  and how well does guided anchoring contain feedback amplification?
  \item \emph{RQ4---Stability and scalability:} Is IRPP-TD stable for near-zero
  residuals or degenerate angular inputs, and how do approximation, convergence,
  runtime, memory, and throughput scale?
  \item \emph{RQ5---Practicality and accountability:} What are the primitive,
  communication, ledger, and end-to-end costs, and how reliably are malicious-DR
  faults corrected under weak challenges and partial audit corruption?
\end{itemize}

\begin{table}[!b]
\centering
\caption{RQ1 Workloads, Vector Semantics, and Realized Participation}
\label{tab:workload-representation}
\scriptsize
\setlength{\tabcolsep}{1.2pt}
\renewcommand{\arraystretch}{1.05}
\begin{tabularx}{\columnwidth}{@{}lXrrc@{}}
\toprule
\textbf{Scene} & \textbf{Task-aligned coordinates} & \(\boldsymbol\ell\) & \textbf{Records} & \(\bar n_{27}/\bar n_{39}\) \\
\midrule
Climate & Pressure, potential temperature, humidity, air density, wind direction & 5 & 420,551 & 26.88/39.42 \\
Traffic & Medium, heavy, micro, long, total, and light-vehicle flow & 6 & 809 & 28.33/38.86 \\
Water & pH, hardness, solids, chloramines, sulfate, conductivity, organic carbon, trihalomethanes, turbidity & 9 & 2,089 & 27.26/39.23 \\
\bottomrule
\end{tabularx}
\vspace{0.2ex}
\parbox{0.98\columnwidth}{\scriptsize Records denotes the clean processed source, not the number of final task reports. Each stored workload contains 100 tasks.}
\end{table}

\textit{Workloads and reports:}
The unchanged legacy pipeline combines stored SUMO trajectories with processed
Climate, Traffic, and Water traces to form six semi-synthetic workloads, each
with 100 tasks in a \(5800\times4300\) coordinate window. Of 100 worker IDs,
1--50 have error probability 0; IDs 51--60, 61--70, 71--80, 81--90, and
91--100 have probabilities 0.1, 0.3, 0.5, 0.7, and 0.9. Target participation
27 and 39 yield the realized means in Table~\ref{tab:workload-representation}
(task ranges 10--56 and 11--67). RQ1 reads the six stored JSON report matrices and
reference vectors directly, without regenerating or injecting test reports.
Each report retains the physical variables naturally associated with its task,
as listed in Table~\ref{tab:workload-representation}; no coordinate replication
or synthetic feature expansion is used. Coordinate-wise
\(Q_{0.95}-Q_{0.05}\) scales are fitted only on tasks 1--20 and pooled across
the two participation settings of the same scene.

\textit{Comparison protocol:}
We compare Mean, CRH, CRH-N, QE, RTD, BLIND, RPPS-TDC, and PRTD through retained
legacy adapters, with byte-identical reports, references, order, and incidence.
The adapters reuse the available legacy code and are not presented as official
implementations; source hashes and adapter decisions are recorded.
Tasks 1--20 freeze \(\zeta=0.25,\theta=0.35\),
\(\mu_1=0.05823,\mu_2=0.001723\), \(s_0=n_b=e_b=5,q_b=2\), and
\(\delta_{\max}=20\); tasks 21--100 are held out and never influence model
selection.

\textit{Metrics and statistics:}
The primary outcomes are native-unit MAE/RMSE within each scene and
scale-normalized NMAE/NRMSE for cross-scene analysis. We additionally record
median/p95 task time, iteration count, across-seed variation, and no-truth rate.
Deterministic baselines are evaluated on all 80 held-out tasks; IRPP-TD uses 30
independent anchor-sampling seeds. Confidence intervals use 2,000 repetitions:
a hierarchical seed/task bootstrap for IRPP-TD and a task bootstrap for
deterministic methods. Paired comparisons resample the same task IDs, while
Table~\ref{tab:main-td-results} reports compact interval half-widths.

\textit{Adversarial and system settings:}
RQ2 uses activity-balanced 50/50 stable populations, controlled cold-start
policies, and forced early rating errors. RQ3 separately varies malicious
ratios, normalized attack strength, collusion, on--off behavior, and reputation
building before attack.
Screening is assessed with F1, FPR/FNR, retention/leakage, anchor purity,
maturity, and recovery delay. RQ5 injects incorrect DR labels, weights, state
transitions, aggregates, payments, and missing/late audit packages while
varying proactive auditing, challenge delay, committee corruption, and
replacement. It reports per-entity computation, communication, ledger growth,
end-to-end latency/throughput, bad-finalization, correction delay, and safe
fallback. The supplement gives attack equations, full grids, audit state
transitions, and implementation versions.

\subsection{Truth-Discovery Effectiveness and Angular Applicability (RQ1)}
\label{sec:rq1-effectiveness}

\begin{table*}[!t]
\centering
\caption{RQ1 Truth-Discovery Accuracy on Held-Out Tasks 21--100 at Target \(\bar n=27\) (Lower Is Better)}
\label{tab:main-td-results}
\footnotesize
\setlength{\tabcolsep}{2.40pt}
\renewcommand{\arraystretch}{0.86}
\newcommand{\dstrut}{\rule[-0.55ex]{0pt}{2.15ex}}
\newcommand{\mci}[2]{\makebox[2.75em][r]{#1}{\scriptsize$\pm$\makebox[1.35em][l]{#2}}}
\begin{tabular}{lcccccc>{\centering\arraybackslash}m{4.80em}>{\centering\arraybackslash}m{4.60em}}
\toprule
Method & \multicolumn{2}{c}{Climate} & \multicolumn{2}{c}{Traffic} & \multicolumn{2}{c}{Water} & \multicolumn{1}{c}{\multirow{2}{*}{\shortstack[c]{Mean\\[-0.35ex]NRMSE}}} & \multicolumn{1}{c}{\multirow{2}{*}{\shortstack[c]{Time/task\\[-0.35ex](ms)}}} \\[-0.42ex]
\cmidrule(lr){2-3}\cmidrule(lr){4-5}\cmidrule(lr){6-7}
 & MAE & RMSE & MAE & RMSE & MAE & RMSE & & \\[-0.30ex]
\midrule
Mean\dstrut & \mci{5.581}{0.61} & \mci{9.288}{1.1} & \mci{16.71}{4.2} & \mci{39.21}{13} & \mci{29.86}{6.1} & \mci{119.23}{28} & 0.0213 & \textbf{4.43e-03} \\
CRH\dstrut & \mci{2.338}{0.29} & \mci{3.798}{0.52} & \mci{6.452}{1.6} & \mci{14.53}{4.0} & \mci{10.18}{3.0} & \mci{48.46}{17} & 0.0116 & \underline{0.2452} \\
CRH-N\dstrut & \mci{0.7762}{0.09} & \mci{1.194}{0.13} & \mci{2.564}{0.70} & \mci{6.622}{2.7} & \mci{6.394}{1.3} & \mci{24.70}{4.8} & \underline{6.47e-03} & 0.5850 \\
QE\dstrut & \mci{5.277}{0.56} & \mci{8.774}{0.95} & \mci{15.75}{3.9} & \mci{36.72}{12} & \mci{28.19}{5.6} & \mci{112.30}{24} & 0.0201 & 4.180 \\
RTD\dstrut & \mci{2.147}{0.16} & \mci{3.459}{0.26} & \mci{6.160}{1.5} & \mci{16.14}{4.5} & \mci{19.10}{4.8} & \mci{82.87}{28} & 0.0237 & 122.17 \\
BLIND\dstrut & \mci{5.281}{0.58} & \mci{8.706}{1.1} & \mci{16.65}{4.1} & \mci{38.50}{9.4} & \mci{24.88}{5.4} & \mci{99.84}{28} & 0.0200 & 20.89 \\
RPPS-TDC\dstrut & \underline{\mci{0.6044}{0.08}} & \underline{\mci{0.8981}{0.12}} & \underline{\mci{2.158}{0.59}} & \underline{\mci{5.035}{1.7}} & \underline{\mci{3.371}{0.70}} & \underline{\mci{11.69}{4.0}} & 9.38e-03 & 0.6059 \\
PRTD\dstrut & \mci{1.762}{0.21} & \mci{2.764}{0.38} & \mci{4.585}{1.1} & \mci{9.681}{2.4} & \mci{6.711}{2.1} & \mci{33.26}{17} & 0.0102 & 1.175 \\
\textbf{IRPP-TD}\dstrut & \textbf{\mci{0.4283}{0.01}} & \textbf{\mci{0.7304}{0.03}} & \textbf{\mci{1.362}{0.08}} & \textbf{\mci{3.660}{0.27}} & \textbf{\mci{2.418}{0.12}} & \textbf{\mci{11.02}{0.74}} & \textbf{1.76e-03} & 1.626 \\[-0.30ex]
\bottomrule
\end{tabular}
\par\vspace{0.70ex}
\noindent\parbox{0.98\textwidth}{\scriptsize
Each \(\pm\) value is the half-width of a 95\% bootstrap interval (hierarchical seed/task for IRPP-TD; task bootstrap for deterministic baselines). IRPP-TD uses 30 anchor-sampling seeds. NMAE/NRMSE use calibration-fitted coordinate scales. Bold and underline mark the best and second-best point estimates under the matched adapter protocol.}
\end{table*}

\begin{figure*}[!t]
\centering
\includegraphics[width=\textwidth]{fig_rq1_main.pdf}
\caption{RQ1 truth-discovery effectiveness on three SUMO-driven workloads.
(a) Held-out NRMSE at \(\bar n=27\) with 95\% bootstrap intervals; colors/markers
denote Climate, Traffic, and Water. (b) Empirical CDF of held-out task NRMSE at
\(\bar n=27\) for IRPP-TD and the three closest retained accuracy competitors.
(c) NRMSE change from target participation 27 to 39. (d) Scene-mean
accuracy--cost operating points. Lower error and cost are better; a left-shifted
CDF indicates lower task error.}
\label{fig:rq1-overview}
\end{figure*}

\textit{Overall accuracy:}
Table~\ref{tab:main-td-results} gives the complete matched comparison on held-out
tasks at \(\bar n=27\). IRPP-TD obtains NRMSEs of 0.00314, 0.000526, and
0.00160 on Climate, Traffic, and Water, respectively, and has the lowest
native-unit MAE and RMSE in every scene. Relative to the strongest retained
normalized baseline, CRH-N, it reduces NRMSE by 47.8\%, 54.1\%, and 87.0\%.
The paired IRPP-TD-minus-CRH-N intervals are
\([-3.49,-2.23]\!\times\!10^{-3}\),
\([-8.43,-4.03]\!\times\!10^{-4}\), and
\([-1.67,-0.33]\!\times\!10^{-2}\), so all three comparisons exclude zero.

The same conclusion is not tied to one baseline family. Against PRTD, IRPP-TD
reduces NRMSE by 76.9\%, 79.9\%, and 89.0\%. Against the closer native-error
competitor RPPS-TDC, its MAE decreases by 29.1\%, 36.9\%, and 28.3\%, while
RMSE decreases by 18.7\%, 27.3\%, and 5.7\%. The smaller Water-RMSE margin and
the visibly wider baseline interval are retained in Table~\ref{tab:main-td-results}
rather than being hidden by rank-only reporting.

\textit{Task-error distribution and participation:}
Fig.~\ref{fig:rq1-overview}(b) retains the 240 held-out scene--task units rather
than reporting only scene means. IRPP-TD's empirical distribution is shifted
toward lower NRMSE and has a shorter high-error tail than CRH-N, PRTD, and
RPPS-TDC. Hence, its mean advantage is not driven by only a small subset of
easy tasks. Supplementary Fig.~S1(d) separates the calibration, early-test, and
mature-test phases. From tasks 21--50 to 51--100, NRMSE falls from 0.00347 to
0.00293 on Climate and from 0.000691 to 0.000396 on Traffic. Water changes from
0.00140 to 0.00170 but remains below its calibration value 0.00200; therefore,
the phase evidence is favorable overall without implying monotone improvement
for every workload.

Fig.~\ref{fig:rq1-overview}(c) further shows that increasing target
participation from 27 to 39 reduces IRPP-TD NRMSE by 21.0\%, 44.0\%, and
20.8\% on the three workloads. IRPP-TD remains the best retained method in all
three scenes at \(\bar n=39\), indicating that its advantage is not confined to
the lower-participation setting. Across 30 anchor-sampling seeds, every held-out
task returns a finite truth; the bounded update averages 3.93 and never exceeds
seven iterations. Across-seed NRMSE standard deviations range from
\(2.32\times10^{-6}\) to \(2.19\times10^{-4}\).

\textit{Accuracy--cost tradeoff:}
Fig.~\ref{fig:rq1-overview}(d) places these gains against per-task runtime.
IRPP-TD averages 1.626~ms per task and a scene-mean NRMSE of
\(1.76\times10^{-3}\). It is not the fastest estimator: CRH-N uses 0.585~ms
and PRTD uses 1.175~ms. However, IRPP-TD lowers their mean NRMSE by 72.9\% and
82.8\%, respectively, and remains substantially faster than BLIND (20.89~ms)
and RTD (122.17~ms). This places IRPP-TD in a favorable millisecond-scale
accuracy--cost region rather than claiming minimum computation in isolation.

\textit{Angular applicability:}
The reports are evaluated in their original task-aligned coordinate systems,
without artificial feature expansion. The three workloads therefore provide
evidence for RABOD when sensing reports contain meaningful, nondegenerate
cross-coordinate geometry. They do not establish an angular advantage for
scalar, rank-deficient, or unpartitioned multi-regime inputs; RQ4 explicitly
tests these boundary cases and the declared bypass/fallback behavior.

\FloatBarrier
\subsection{Cold-Start Efficiency and Recovery (RQ2)}
\label{sec:rq2-feedback}

\textit{Protocol and population construction:}
RQ2 preserves every RQ1 truth vector, task order, and SUMO-determined
participation ID, but regenerates reports to expose worker-level ground truth.
IDs are activity-stratified and interleaved before assigning exactly 50 stable
high-quality (H) and 50 stable low-quality (L) workers, decoupling type from ID
and mobility frequency. H reports use truncated $\pm1\%$ noise; L reports
corrupt every coordinate by a signed $5\%$--$20\%$ deviation. Worker Macro-F1
penalizes immature/uncertain states as abstentions, and NRMSE measures
end-to-end accuracy.

IRPP-HQ5 is the submitted policy: while fewer than $n_b=5$ mature ordinary
anchors are present, it adds at most five disjoint high-quality reports.
IRPP-HQ10 doubles that cap. ABOD-HQ5 changes only the angular evaluator: it
uses every reputation-guided anchor instead of RABOD sampling and receives the
same five-report bootstrap. No-Extra adds no report and uses a temporary
all-ordinary reference only when fewer than four qualified anchors exist;
Random-5 adds five reports from a disjoint half-H/half-L pool. HQ5 is a nested
prefix of HQ10. All policies share reports, normalization, reputation state,
stopping rules, seeds, and tuning budget. Because exact-ABOD scores have a
different scale, its two cutoffs are selected on RQ1 tasks 1--20 with the same
$3\times3$ quantile grid at the frozen $\theta$, never on RQ2 outcomes.

All reputation and RABOD thresholds remain frozen from RQ1. Under E0, natural
maturation is measured from task start. E1 forces a false-low update at an H
worker's second participation; E3 does so at participations 2--4. Stable
completion requires three consecutive high classifications; unresolved E1/E3
workers are right-censored. We report ten-task blocks and 95\% intervals from
2,000 bootstraps of 30 paired scene-macro seeds; $\bar n=39$ replication and
resource details are in the supplement.

\begin{figure*}[!t]
\centering
\includegraphics[width=\textwidth]{fig_rq2_coldstart_main_recovery_1x4.pdf}
\caption{RQ2 cold-start efficiency and recovery at $\bar n=27$.
(a) Abstention-aware Worker Macro-F1; the inset enlarges blocks 30--60.
(b) Mature NRMSE versus angular-pair evaluations per task; marker area gives
the external-report cap (0, 5, or 10) per active task. (c) Paired block
$\Delta$NRMSE under E1/E3 relative to E0. (d) Stable completion: E0 is measured
from task start; E1/E3 after the final forced error, with unresolved workers
censored. Curves/bands are scene-macro means/95\% intervals over 30 paired
seeds. Higher is better in (a,d), lower in (b), and zero is ideal in (c).}
\label{fig:rq2-coldstart}
\end{figure*}

\begin{table*}[!t]
\centering
\caption{RQ2 Component Ablation of Screening, Anchor Construction, TD
Accuracy, and Bootstrap Cost at Target $\bar n=27$}
\label{tab:rq2-main}
\scriptsize
\setlength{\tabcolsep}{3.25pt}
\renewcommand{\arraystretch}{0.94}
\begin{tabular}{lccccccc}
\toprule
Variant & AUPRC$\uparrow$ & Screening-F1$\uparrow$ & Anchor purity$\uparrow$ &
Clean retention$\uparrow$ & Bad leakage$\downarrow$ &
NRMSE ($10^{-3}$)$\downarrow$ & Seed/task$\downarrow$ \\
\midrule
\textbf{Full (IRPP-TD)} & 0.994$\pm$0.002 & 0.972$\pm$0.006 & 0.959$\pm$0.010 & 0.951$\pm$0.005 & 0.009$\pm$0.003 & 1.71$\pm$0.23 & 0.86$\pm$0.03 \\
No-U & 0.979$\pm$0.001 & 0.937$\pm$0.005 & 0.948$\pm$0.008 & 0.916$\pm$0.008 & 0.059$\pm$0.006 & 1.68$\pm$0.21 & 0.86$\pm$0.03 \\
Binary-Rep. & 0.984$\pm$0.001 & 0.963$\pm$0.006 & 0.942$\pm$0.009 & 0.928$\pm$0.014 & 0.021$\pm$0.003 & 1.74$\pm$0.22 & 0.93$\pm$0.06 \\
All-Anchors & 0.552$\pm$0.008 & 0.196$\pm$0.012 & 0.507$\pm$0.001 & 0.920$\pm$0.013 & 0.682$\pm$0.009 & 12.15$\pm$1.30 & 0.76$\pm$0.03 \\
Sequential & 0.552$\pm$0.008 & 0.196$\pm$0.012 & 0.507$\pm$0.001 & 0.894$\pm$0.016 & 0.608$\pm$0.012 & 11.56$\pm$1.20 & 0.76$\pm$0.03 \\
No-Seeds & -- & 0.000$\pm$0.000 & -- & 1.000$\pm$0.000 & 1.000$\pm$0.000 & 13.62$\pm$1.50 & 0.00$\pm$0.00 \\
Always-Seeds & 0.995$\pm$0.001 & 0.977$\pm$0.005 & 0.960$\pm$0.010 & 0.953$\pm$0.005 & 0.008$\pm$0.002 & 1.39$\pm$0.17 & 5.00$\pm$0.00 \\
\bottomrule
\end{tabular}
\par\vspace{0.45ex}
\noindent\parbox{0.98\textwidth}{\scriptsize
Values are scene-macro means $\pm$ the larger 95\% hierarchical-bootstrap
half-width over 30 paired seeds. Anchor purity treats the zero-error ordinary
registry as reliable and excludes external seeds. AUPRC is computed only where
angular evidence exists; '--' is undefined, not zero. Full combines ternary
TV-uncertainty reputation, reputation-guided anchors, and adaptive seeds.
No-U removes uncertainty shrinkage; Binary-Rep. replaces the ternary state by
a binary Beta state; All-Anchors forms the angular reference from all current
workers; Sequential applies all-worker angular detection before reputation
gating; No-Seeds disables bootstrap; Always-Seeds keeps all five seeds active.}
\end{table*}

\textit{Cold-start accuracy and cost:}
By task block 30, IRPP-HQ5 and IRPP-HQ10 reach Macro-F1 0.898 and 0.904,
compared with 0.851 for ABOD-HQ5 and 0.269 and 0.254 for No-Extra and Random-5
(Fig.~\ref{fig:rq2-coldstart}(a)). Their final values are 0.99994, 1.000,
0.99788, 0.358, and 0.367, respectively. Mature NRMSE is 0.001639 for HQ5,
0.001633 for HQ10, and 0.001730 for exact ABOD, versus 0.02460 and 0.02385 for the
two controls. HQ5 evaluates 430 angular pairs per task, 81.2\% fewer than
exact ABOD's 2,292. Doubling the trusted cap improves mature NRMSE by only
$6.22\times10^{-6}$ while increasing cumulative external reports from 90.4 to
178.3. Thus, five trusted reports give the stronger accuracy--worker-cost
operating point; random extra reports do not substitute for trusted geometry.

\textit{Finite early-error recovery:}
E1 changes full-horizon NRMSE by $6.1\times10^{-6}$ and final Macro-F1 by
$-0.04$ percentage points. E3 increases NRMSE by $8.54\times10^{-5}$, peaks at
$2.93\times10^{-4}$ in block 20, and leaves a 1.07-point final F1 loss. In
Fig.~\ref{fig:rq2-coldstart}(d), E0 naturally matures after a median of 7 own
participations with no censoring. Measured instead from the final forced error,
median recovery is 6 participations for E1 and 10 for E3; restricted means are
7.26 and 11.26, with 0.24\% and 4.96\% censoring. E0 is therefore a natural
maturation reference, not a same-origin recovery curve. The feedback loop
recovers from a limited early mistake, while repeated mistakes produce a
visible and longer tail rather than an irreversible lockout.

\textit{Component isolation:}
Table~\ref{tab:rq2-main} retains the RQ1 graded workers, isolating seven
mechanisms from the binary stress population.
Full reaches 0.994 AUPRC, 0.972 screening-F1, 0.959 anchor purity, and NRMSE
$1.71\times10^{-3}$. No-U yields 0.979 AUPRC, 0.937 F1, 0.948 purity, and
0.059 bad leakage; Binary-Rep. yields 0.984, 0.963, 0.942, and 0.021.
Their NRMSEs remain near Full at $1.68\times10^{-3}$ and $1.74\times10^{-3}$,
so uncertainty and the ternary state
chiefly improve screening/anchor point estimates here. All-Anchors and
Sequential reduce purity to 0.507 and raise NRMSE to 0.01215 and 0.01156, showing
that reputation must shape the angular reference. No-Seeds deadlocks under the
prior; Always-Seeds lowers NRMSE to 0.00139 at five reports per task. Together,
the table and Fig.~\ref{fig:rq2-coldstart} separate component necessity,
cold-start benefit, approximation cost, and finite-error recovery.

\FloatBarrier
\subsection{Robustness to Collusion and Reputation Poisoning (RQ3)}
\label{sec:rq3-robustness}

\begin{figure*}[!t]
\centering
\includegraphics[width=\textwidth]{fig_rq3_main_1x4.pdf}
\caption{RQ3 worker-side robustness at $\bar n=27$. (a) Compact-collusion
NRMSE versus prevalence. (b) Active-task NRMSE across clean and four attacks.
(c) Malicious-report leakage under method-faithful hard screening. (d) Ordinary-anchor
purity after the task-41 switch. RPPS-TDC uses native $p=0.3$; PRTD's added
gate targets 96\% clean acceptance. Neither threshold affects (a,b).
Attacks in (b,c) use $(\rho_m,\kappa)=(0.3,0.5)$. Estimates carry 95\%
paired-seed intervals over 30 seeds; lower is better in (a)--(c), higher in (d).}
\label{fig:worker-robustness}
\end{figure*}

\textit{Stress protocol:}
We reuse the frozen RQ1 methods and $\bar n=27$ SUMO incidence. Activity-balanced
sets vary \(\rho_m=0{:}0.1{:}0.8\); honest reports stay within $\pm1\%$, and
attack displacement has calibration-scale RMS \(\kappa\). Fig.~\ref{fig:worker-robustness}(c) uses native
IRPP-TD/RPPS-TDC decisions and a PRTD gate fitted only on clean tasks 1--20;
all precede attack evaluation. HQ5 remain honest; details are in the supplement.

\textit{Accuracy across attacks:}
Compact collusion raises IRPP-TD from 0.00113 NRMSE to only 0.00215 at
\(\rho_m=0.8\), versus the closest baseline at 0.406 and no abstention
(Fig.~\ref{fig:worker-robustness}(a)). In panel (b), IRPP-TD records 0.00137
under both immediate attacks; the nearest Independent and Compact baselines
give 0.00886 and 0.04045. Its clean 0.00113 is close to QE's 0.00107.
Under On--Off, IRPP-TD/CRH-N give 0.03631 and 0.03544; under Mature-Anchor,
IRPP-TD is lowest at 0.03315. It therefore uniquely preserves near-clean
accuracy across both immediate geometries and remains competitive when timed.

\textit{Leakage and feedback:}
Across Independent, Compact, On--Off, and Mature-Anchor, Fig.~\ref{fig:worker-robustness}(c) gives
IRPP-TD leakage of 0, 0, 0.249, and 0.206; RPPS-TDC gives 0.00025, 0.080,
0.293, and 0.295; PRTD retains 0.719--0.780. IRPP-TD accepts 0.953--0.993 of
honest attack-phase reports, above RPPS-TDC's 0.730--0.986; PRTD exceeds
0.999 but admits most poison. Thus, IRPP-TD has the strongest joint leakage--
collateral-rejection tradeoff. Mature no-truth is only 0.019\%. In panel (d),
Full purity recovers from 0.700 to 0.968 and its error ratio to 1.35, whereas
unguided variants stay near 0.697. IRPP-TD therefore combines immediate
accuracy and rejection with lower timed leakage and guided recovery.

\FloatBarrier
\subsection{Numerical Stability and Algorithmic Scalability (RQ4)}
\label{sec:rq4-scalability}

\begin{figure*}[!b]
\centering
\panelbox{(a) Runtime and slope}{Worker evaluation, RABOD, and TD runtime versus \(n_T\); fit the empirical complexity slope after the anchor cap.}
\hfill
\panelbox{(b) Memory and throughput}{Peak RAM and reports/s or tasks/s versus \(n_T\), report dimension, and task horizon.}
\hfill
\panelbox{(c) Dimension, cap, and fidelity}{Vary \(\ell\) and \(\delta_{\max}\); report exact-ABOD score Spearman, filtered-set Jaccard, and F1 gap.}
\hfill
\panelbox{(d) Convergence and degeneracy}{Iterations/NRMSE versus \(\epsilon,t_{\max}\); summarize zero residual, duplicate, low-rank, scalar, and empty-set fallback.}
\caption{Algorithmic runtime, memory/throughput, approximation fidelity,
convergence, and numerical fallback behavior (RQ4).}
\label{fig:algorithm-scalability}
\end{figure*}

Fig.~\ref{fig:algorithm-scalability}(a)--(c) varies \(n_T\), dimension, and
\(\delta_{\max}\), reporting module time, peak RAM, throughput, empirical
complexity slope, and agreement with exact ABOD (score Spearman, retained-set
Jaccard, and F1 gap). At \(n_T=10^4\), RABOD uses \tbd{} s/\tbd{} MB with
slope \tbd{} and agreement \tbd; this slope applies only to the angular stage.

Panel (d) varies \(\epsilon,t_{\max}\) and tests zero residuals, identical,
duplicate/collinear, low-rank, scalar, insufficient-pair, and empty-retained-set
inputs. Across \tbd{} trials it observes \tbd{} NaN/Inf values and \tbd{} cap
hits. The main text retains the \((n_T,\delta_{\max})\) and
\((\rho_m,\theta)\) interactions; all other sensitivity and fallback grids are
reported in the supplement.

\FloatBarrier
\subsection{End-to-End Overhead and Conditional Accountability (RQ5)}
\label{sec:rq5-system}

\begin{table*}[!t]
\centering
\caption{End-to-End Comparison under a Matched MCS Workload and Security Setting}
\label{tab:e2e-matched-comparison}
\scriptsize
\setlength{\tabcolsep}{1.7pt}
\renewcommand{\arraystretch}{1.04}
\begin{tabularx}{\textwidth}{@{}>{\raggedright\arraybackslash}p{0.08\textwidth}>{\raggedright\arraybackslash}p{0.16\textwidth}>{\raggedright\arraybackslash}p{0.23\textwidth}>{\raggedright\arraybackslash}p{0.20\textwidth}ccc@{}}
\toprule
\textbf{Scheme} & \textbf{Analytics} & \textbf{Privacy scope} &
\textbf{Malicious-DR handling} & \textbf{Latency} &
\textbf{Traffic/report} & \textbf{Chain storage/task} \\
\midrule
PRTD & Reputation-weighted TD & \tbd{} from exact threat model & \tbd{} from exact protocol & \tbd & \tbd & \tbd{} / N/A \\
RPPS-TDC & Reputation-aware collection/TD & \tbd{} from exact threat model & \tbd{} from exact protocol & \tbd & \tbd & \tbd \\
BSIF & Task evaluation/incentive; no TD & \tbd{} from exact threat model & \tbd{} from exact protocol & \tbd & \tbd & \tbd \\
\textbf{IRPP-TD} & RABOD and bounded TD & TP/RB/public; excludes intended DR & Proactive-or-challenged conditional accountability & \textbf{\tbd} & \textbf{\tbd} & \textbf{\tbd} \\
\bottomrule
\end{tabularx}
\vspace{0.25ex}
\parbox{0.98\textwidth}{\scriptsize
\textit{Fairness rule:} use the same worker count, payload, hardware, network,
and security level. Compare only shared protocol phases. BSIF has no TD module,
so its TD accuracy is N/A. Never mix locally measured and source-reported
runtime values in one quantitative ranking.}
\end{table*}

\begin{table}[!t]
\centering
\caption{Malicious-DR Fault Injection and Expected Protocol Path}
\label{tab:dr-fault-matrix}
\scriptsize
\setlength{\tabcolsep}{1.35pt}
\renewcommand{\arraystretch}{1.05}
\begin{tabularx}{\columnwidth}{@{}>{\raggedright\arraybackslash}p{0.26\columnwidth}>{\raggedright\arraybackslash}p{0.31\columnwidth}X@{}}
\toprule
\textbf{Injected deviation} & \textbf{Trigger condition} & \textbf{Expected path/result} \\
\midrule
Wrong label/weight/state & Challenge or proactive audit & Re-execute; correct or reject pending batch \\
Biased aggregate/payment & Challenge or proactive audit & Correct aggregate/payment; slash on DR fault \\
Missing/late package & Audit deadline expires & Default DR fault and correction path \\
First committee unavailable & Fewer than \(t_A\) timely votes & Draw one replacement committee \\
Second committee unavailable & Replacement also fails & Preserve predecessor; release safe base payment \\
Threshold corruption & At least \(t_A\) corrupt auditors & Declared residual failure boundary \\
\bottomrule
\end{tabularx}
\end{table}

\begin{figure*}[!t]
\centering
\panelbox{(a) Primitive/entity time}{KEM, AEAD, signatures, commitments, and Merkle operations on worker, TP, DR, RB, and auditor hardware; report median/p95.}
\hfill
\panelbox{(b) P1--P7 system cost}{Per-protocol latency, traffic, storage/task, TPS, and confirmation delay versus workers and payload length.}
\hfill
\panelbox{(c) Matched end-to-end systems}{IRPP-TD, BSIF, PRTD, and RPPS-TDC common-path latency/traffic on one platform and security level.}
\hfill
\panelbox{(d) Accountability Pareto}{Audit cost versus bad-finalization and correction delay under no/delayed challenge and partial committee corruption.}
\caption{Concrete primitive and protocol costs, matched end-to-end system
comparison, and the conditional-accountability tradeoff (RQ5).}
\label{fig:system-accountability}
\end{figure*}

The concrete suite is \fillnote{insert KEM/AEAD/signature/hash parameters and
versions}; the FISCO BCOS testbed uses \tbd{} nodes, \tbd{} consensus, \tbd{}
RTT, and \tbd-byte reports. Fig.~\ref{fig:system-accountability}(a)--(c) reports
per-entity median/p95 time, per-protocol traffic/storage, TPS, confirmation
delay, and matched common-path costs. The normal/audited paths require
\tbd/\tbd{} ms and \tbd/\tbd{} bytes; Table~\ref{tab:e2e-matched-comparison}
keeps analytic scope and privacy boundaries visible beside these costs.

Table~\ref{tab:dr-fault-matrix} preregisters the correct ledger path for each
fault. Panel (d) varies no/timely/delayed challenge, \(p_A\), committee
corruption, timeout, and one replacement, reporting fault detection, bad
finalization, correction delay, and safe fallback. It obtains \tbd{} bad
finalization at \tbd{} extra cost; fewer than \(t_A\) corrupt auditors correct
\tbd{} of triggered faults. For observed timely-challenge probability \(c_T\),
we report \(q_T=1-(1-p_A)(1-c_T)\) and
\(\bar C_T=(1-q_T)C_N+q_TC_A\) beside empirical outcomes. Threshold corruption
and active RB--DR collusion remain failure boundaries, not solved cases.

% Keep the concluding validity discussion with the final experiment page;
% all RQ5 floats have already been introduced above and may settle naturally.
\subsection{Summary and Threats to Validity}
\label{sec:experiment-summary}

RQ1 shows that IRPP-TD improves over every retained baseline, benefits from
higher participation, returns a finite truth on every held-out run, and uses at
most seven iterations on the tested representations. RQ2 shows that an adaptive
five-report cap resolves the activity-balanced 50/50 cold start, matches or
improves exact-ABOD accuracy with 81.2\% fewer angular pairs, and recovers from
limited early false-low updates. The ablation identifies necessary
anchor/bootstrap components; uncertainty and ternary state improve
screening/anchor point estimates without a resolved TD-error gain. RQ3 shows
strong suppression of immediate poisoning and eventual guided-anchor recovery,
but also exposes transient errors under reputation-aware timing. Conclusions
for RQ4--RQ5 remain \tbd{} until those experiments are executed.

\textit{Validity and reproducibility:}
Real reference vectors are paired with semi-synthetic reliability/attacks; the
three domains do not cover scalar, multi-regime, or out-of-model collusion.
Methods share inputs and disjoint splits; test truth is evaluation-only, and
seeded sequences are the bootstrap units. Complete artifacts are released;
unmatched systems remain N/A or \tbd{}.

\end{document}
